
\scanpage{63}
\mspage{19}
La question essentielle est donc comment reconstituer, pour un graphe \deleted{$\unsure{\Gamma}$} pondéré à
\deleted{à repères donné} \inserted{repères} (disons) $\underline\Gamma$, le
$\mathcal T_{\mathrm{top}}^\pm(\underline\Gamma)$,
\deleted{plus} \inserted{avec} sa dépendance fonctorielle en $\underline\Gamma$,
\deleted{\unsure{des foncteurs des}} \inserted{et l'explicitation}
des sept opérations, à l'exception de $4^\circ)$ et $5^\circ)$.
\margin{\uline{NB} on laisse ici tomber le ``$S$'' ou ``$D$'', contenus déjà dans $\underline\Gamma$.}

Pour l'opération $4^\circ)$, elle sera obtenue automatiquement, par la description explicite des
$\mathcal T_{\mathrm{top}}^\pm(\underline\Gamma)$
(qui sera faite en $=$ temps et de la $=$ façon que pour les $\unsure{S}\mathcal T_{g,\nu}^\pm$ et $\mathcal T_{g,\nu}^\pm$ \unsure{ci-dessus}),
et le bouchage des trous \inserted{pourrait} s'expliciter à l'aide \deleted{\unsure{des 7 opérations}} \inserted{des cas} particuliers de ces opérations pour les
\deleted{analogues des}
$S\mathcal T_{g,\nu}$ et $\mathcal T_{g,\nu}$, qu'il faudrait jeter dans le sac des
\deleted{\unsure{données}} \inserted{structures} ``élémentaires'' à déterminer.

\deleted{$(**)$}\quad
\diagram{\begin{tikzcd}[ampersand replacement=\&,column sep=huge]
\deleted{S}\mathcal T_{g,\nu,\nu'}^\pm\arrow[r]
\&\deleted{S}\mathcal T_{g,\nu-\nu'}
\end{tikzcd}}
\margin{avec ou sans $S$?}
\inserted{les $X\in\operatorname{Ob}\mathcal T_{g,\nu}^\pm$, avec la donnée d'un sous ens. $I'\subset I(X)$ avec $\operatorname{card}I'=\nu'$}
\mspage{62}
